%================================================================
\chapter{Results and Discussion}
\label{ch:results}
%================================================================

\section{Overview}

This chapter presents the experimental results across four domains:
preprocessing and dataset characterisation (\S\ref{sec:preprocess-results}),
model performance on the sentiment classification task
(\S\ref{sec:classification-results}), temporal predictive analysis
(\S\ref{sec:temporal-results}), and Explainable Explainable AI audit findings
(\S\ref{sec:xai-results}). All reported metrics are computed on the
held-out 20\% test partition (319,628 instances) unless otherwise stated.
Figures and tables are generated directly from the experimental outputs
of the implemented pipeline described in Chapter~\ref{ch:impl}.

%================================================================
\section{Preprocessing and Dataset Characterisation}
\label{sec:preprocess-results}
%================================================================

The preprocessing pipeline produced a cleaned corpus of
\textbf{1,598,142} usable tweet instances after deduplication and
removal of empty post-clean rows.
Figure~\ref{fig:eda-distribution} illustrates the class distribution
and tweet length profile of the full corpus.

\begin{figure}[H]
  \centering
  \includegraphics[width=0.95\textwidth]{figures/01_eda_distribution.png}
  \caption{Class distribution (bar and pie charts) and tweet character
           length distribution for the Sentiment140 corpus.}
  \label{fig:eda-distribution}
\end{figure}

The class distribution is substantially imbalanced: positive tweets
account for \textbf{49.98\%} of the corpus, negative tweets
\textbf{49.92\%}, and neutral tweets only \textbf{0.10\%} --- a direct
consequence of the distant supervision annotation methodology
\citep{go2009twitter}. The median tweet length is \textbf{74}
characters, with a mean of \textbf{73.6} characters, reflecting the
characteristically concise nature of Twitter discourse. Figure~\ref{fig:tweet-volume}
shows the daily tweet volume across the corpus's temporal coverage.

\begin{figure}[H]
  \centering
  \includegraphics[width=0.95\textwidth]{figures/01_tweet_volume.png}
  \caption{Daily tweet volume across the Sentiment140 corpus temporal
           coverage (April--June 2009), showing weekly periodic patterns.}
  \label{fig:tweet-volume}
\end{figure}

Figure~\ref{fig:vader-dist} shows the distribution of \gls{vader}
compound scores across the three sentiment classes, confirming that
\gls{vader} is a valid discriminative feature. The mean compound
score for positive-labelled tweets was $\bar{s}_{+} = 0.312$
(SD $= 0.289$), for negative-labelled tweets
$\bar{s}_{-} = -0.267$ (SD $= 0.276$), and for neutral-labelled
tweets $\bar{s}_{0} = 0.041$ (SD $= 0.198$). The clear separation
between positive and negative distributions validates the inclusion
of \gls{vader} compound score as a standalone feature alongside the
\gls{tfidf} representation.

\begin{figure}[H]
  \centering
  \includegraphics[width=0.95\textwidth]{figures/02_vader_distributions.png}
  \caption{\gls{vader} compound score distributions per sentiment class
           (training set). Dashed vertical lines indicate class means.}
  \label{fig:vader-dist}
\end{figure}

\gls{smote} application to the 10\% training subsample balanced the
class distribution from the original ratio of approximately
50\%:0.1\%:50\% (positive:neutral:negative) to an equal distribution
of approximately 33.3\% per class, producing a balanced training set
of \textbf{127,965} instances (\textbf{42,655} per class after
oversampling), as illustrated in Figure~\ref{fig:smote-balance}.

\begin{figure}[H]
  \centering
  \includegraphics[width=0.75\textwidth]{figures/02_smote_balance.png}
  \caption{Class distribution before and after \gls{smote} oversampling
           applied to the 10\% training subsample. The neutral class was
           increased from 128 to 42,655 instances.}
  \label{fig:smote-balance}
\end{figure}

%================================================================
\section{Classification Performance Results}
\label{sec:classification-results}
%================================================================

\subsection{Standard Classification Metrics}
\label{subsec:standard-results}

Table~\ref{tab:standard-metrics} presents the classification performance
metrics for all four models on the held-out test set. DistilBERT was not
available in this experimental run due to the absence of a \gls{gpu}
environment during execution; the \gls{mlp} Neural Network therefore
represents the highest-performing model. The \gls{mlp} achieves the
highest accuracy of \textbf{82.4\%}, consistent with the expected
performance range for feedforward neural network architectures on
the Sentiment140 task \citep{hartmann2022feeling}.

\begin{table}[H]
  \centering
  \caption{Standard classification performance metrics for all four models
           (weighted averages). Bold indicates best performance per column.}
  \label{tab:standard-metrics}
  \begin{tabular}{@{}lcccc@{}}
    \toprule
    \rowcolor{uwsblue!90}
    \textcolor{white}{\textbf{Model}}
      & \textcolor{white}{\textbf{Acc.}}
      & \textcolor{white}{\textbf{Prec.\ (W)}}
      & \textcolor{white}{\textbf{Rec.\ (W)}}
      & \textcolor{white}{\textbf{F1 (W)}} \\
    \midrule
    \rowcolor{tablegrey}
    Logistic Regression  & 0.782 & 0.795 & 0.782 & 0.788 \\
    Random Forest        & 0.813 & 0.825 & 0.813 & 0.818 \\
    \rowcolor{tablegrey}
    MLP Neural Network   & \textbf{0.824} & \textbf{0.835} & \textbf{0.824} & \textbf{0.829} \\
    Ensemble (LR+RF)     & 0.820 & 0.831 & 0.820 & 0.825 \\
    \bottomrule
  \end{tabular}
  \\[4pt]
  \small\textit{Note: Acc.\ = Accuracy; W = weighted average; Prec.\ = Precision; Rec.\ = Recall. Bold indicates best performance.
  DistilBERT excluded: GPU not available in this execution run.}
\end{table}

The performance gap between Logistic Regression and the neural models
(approximately \textbf{4.2} percentage points in accuracy) is consistent
with the broader literature on classical versus ensemble methods for
text classification, reflecting the capacity of ensemble and non-linear
models to capture higher-order feature interactions that linear
classifiers cannot represent \citep{breiman2001random}.

Figure~\ref{fig:rf-importance} shows the top-20 most important features
identified by the Random Forest classifier. The \gls{vader} compound
score ranks as the single most influential feature, confirming the
methodological value of the hybrid feature engineering approach.

\begin{figure}[H]
  \centering
  \includegraphics[width=0.85\textwidth]{figures/03_rf_feature_importances.png}
  \caption{Top-20 feature importances from the trained Random Forest
           classifier. \gls{vader} compound score is the most influential
           single feature, validating the hybrid feature engineering strategy.}
  \label{fig:rf-importance}
\end{figure}

Figure~\ref{fig:mlp-loss} shows the training and validation loss curve
of the \gls{mlp} neural network, demonstrating smooth convergence with
early stopping triggered at epoch \textbf{14} (final training loss:
$0.394$; final validation loss: $0.497$).

\begin{figure}[H]
  \centering
  \includegraphics[width=0.75\textwidth]{figures/03_mlp_loss_curve.png}
  \caption{\gls{mlp} training and validation loss curves. Early stopping
           was triggered at epoch 14, preventing overfitting while
           maintaining generalisation performance.}
  \label{fig:mlp-loss}
\end{figure}

The normalised confusion matrices for all four models are shown in
Figure~\ref{fig:conf-matrices}. A consistent pattern is apparent: the
Neutral sentiment class has the lowest per-class recall across all
models, confirming that neutral sentiment is the most challenging
category to classify correctly. This finding directly validates the
argument by \citet{tang2014neutral} that neutral sentiment occupies
a semantically ambiguous boundary position.

\begin{figure}[H]
  \centering
  \includegraphics[width=\textwidth]{figures/05_confusion_matrices.png}
  \caption{Normalised confusion matrices for all four models on the
           held-out test set. Values represent true positive rates per
           class. The Neutral class consistently shows the lowest recall.}
  \label{fig:conf-matrices}
\end{figure}

\subsection{ROC-AUC Analysis}
\label{subsec:roc-results}

Table~\ref{tab:roc-auc} presents the multi-class \gls{roc}-\gls{auc}
results computed using the \gls{ovr} strategy. Figure~\ref{fig:roc-curves}
shows the corresponding \gls{roc} curve plots.

\begin{table}[H]
  \centering
  \caption{Multi-class \gls{roc}-\gls{auc} results (One-vs-Rest) and
           training time for all four models.}
  \label{tab:roc-auc}
  \begin{tabular}{@{}lcccccr@{}}
    \toprule
    \rowcolor{uwsblue!90}
    \textcolor{white}{\textbf{Model}}
      & \textcolor{white}{\textbf{AUC Neg}}
      & \textcolor{white}{\textbf{AUC Neu}}
      & \textcolor{white}{\textbf{AUC Pos}}
      & \textcolor{white}{\textbf{Macro AUC}}
      & \textcolor{white}{\textbf{Wt.\ AUC}}
      & \textcolor{white}{\textbf{Train Time}} \\
    \midrule
    \rowcolor{tablegrey}
    Logistic Regression & 0.9145 & 0.9576 & 0.9145 & 0.9289 & 0.9146 & 4m 7s \\
    Random Forest       & 0.9307 & 0.9703 & 0.9306 & 0.9439 & 0.9307 & 18m 7s \\
    \rowcolor{tablegrey}
    MLP Neural Network  & \textbf{0.9362} & 0.9706 & \textbf{0.9361} & 0.9477 & \textbf{0.9362} & 27m 14s \\
    Ensemble (LR+RF)    & 0.9341 & \textbf{0.9730} & 0.9343 & \textbf{0.9471} & 0.9342 & 5m 12s \\
    \bottomrule
  \end{tabular}
  \\[4pt]
  \small\textit{Note: Wt.\ = Weighted average. Bold = best per column.}
\end{table}

\begin{figure}[H]
  \centering
  \includegraphics[width=\textwidth]{figures/05_roc_curves.png}
  \caption{\gls{roc}-\gls{auc} curves for all four models using the
           One-vs-Rest strategy. Dashed diagonal line represents the
           random classifier baseline (\gls{auc} $= 0.500$).}
  \label{fig:roc-curves}
\end{figure}

Several important findings emerge from the \gls{roc}-\gls{auc} analysis.
First, all models achieve \gls{auc} values substantially above the random
classifier baseline of 0.5 across all three sentiment classes, confirming
genuine discriminative ability for every architecture. Second, contrary to
the initial hypothesis, the Neutral class \gls{auc} is \emph{higher} than
the Positive and Negative class \gls{auc} values across all models --- for
example, Logistic Regression achieves Neutral \gls{auc} = 0.9576 versus
Negative \gls{auc} = 0.9145. This result can be attributed to the
\gls{ovr} strategy: with so few neutral instances (0.10\% of the corpus),
the model rarely predicts neutral and therefore rarely makes neutral-class
false positives, producing inflated \gls{auc} scores for that class. This
phenomenon is a documented artefact of \gls{ovr} \gls{auc} computation
on highly imbalanced classes \citep{hand2001simple}, and underscores the
importance of examining per-class performance in the confusion matrix
alongside \gls{auc} scores.

The Ensemble (LR+RF) achieves the highest macro \gls{auc} of
\textbf{0.9471} while also achieving the highest Neutral class \gls{auc}
of \textbf{0.9730}, representing the best balance between performance and
computational cost (training time: 5m 12s). The \gls{mlp} achieves
marginally higher weighted \gls{auc} (0.9362 vs 0.9342) at the cost of
27 minutes of training time, suggesting that the Ensemble is the optimal
choice for production deployment in resource-constrained settings.

%================================================================
\section{Temporal Sentiment Analysis Results}
\label{sec:temporal-results}
%================================================================

Figure~\ref{fig:temporal-dynamics} presents the four-panel temporal
dynamics plot across the 81-day observation window of the Sentiment140
corpus (6 April to 25 June 2009).

\begin{figure}[H]
  \centering
  \includegraphics[width=\textwidth]{figures/04_temporal_dynamics.png}
  \caption{Temporal sentiment dynamics: (a) daily mean \gls{vader}
           compound score with 7-day rolling mean smoothing, (b) daily
           sentiment velocity, (c) daily tweet volume with 7-day rolling
           mean, and (d) sentiment class proportion stacked area chart.}
  \label{fig:temporal-dynamics}
\end{figure}

The daily mean compound \gls{vader} score ranged from \textbf{0.051}
to \textbf{0.184} across the observation period, with a mean of
\textbf{0.126} (SD $= 0.028$), indicating a consistent modest net
positive sentiment bias in the corpus. Visual inspection of the
compound score time series confirms periodic weekly cycles --- higher
positive sentiment at weekends relative to weekdays --- justifying the
application of the 7-day rolling mean smoothing prior to velocity
computation. Sentiment velocity showed a mean absolute value of
\textbf{0.004} (SD $= 0.003$).

Figure~\ref{fig:lag-corr} and Table~\ref{tab:lag-corr} present the
Pearson lag correlation results between sentiment velocity at time $t$
and tweet volume at time $t+k$ for $k = 1, \ldots, 7$ days.

\begin{figure}[H]
  \centering
  \includegraphics[width=0.80\textwidth]{figures/04_lag_correlation.png}
  \caption{Pearson correlation coefficients between sentiment velocity
           and next-$k$-day tweet volume ($k = 1$--$7$ days). Green bars
           indicate statistical significance ($p < 0.05$); grey bars are
           not significant.}
  \label{fig:lag-corr}
\end{figure}

\begin{table}[H]
  \centering
  \caption{Lag correlation results: Pearson $r$ between sentiment velocity
           at time $t$ and tweet volume at time $t+k$.}
  \label{tab:lag-corr}
  \begin{tabular}{@{}cccc@{}}
    \toprule
    \rowcolor{uwsblue!90}
    \textcolor{white}{\textbf{Lag $k$ (days)}}
      & \textcolor{white}{\textbf{Pearson $r$}}
      & \textcolor{white}{\textbf{$p$-value}}
      & \textcolor{white}{\textbf{Significant?}} \\
    \midrule
    \rowcolor{tablegrey} 1 & 0.412 & 0.0031 & $\checkmark$ \\
                          2 & 0.293 & 0.0278 & $\checkmark$ \\
    \rowcolor{tablegrey} 3 & 0.181 & 0.1402 & $\times$ \\
                          4 & 0.124 & 0.2891 & $\times$ \\
    \rowcolor{tablegrey} 5 & 0.089 & 0.4712 & $\times$ \\
                          6 & 0.061 & 0.6183 & $\times$ \\
    \rowcolor{tablegrey} 7 & 0.042 & 0.7294 & $\times$ \\
    \bottomrule
  \end{tabular}
  \\[4pt]
  \small\textit{Significance threshold $\alpha = 0.05$.
  $\checkmark$ = significant; $\times$ = not significant.}
\end{table}

A statistically significant positive Pearson correlation of
$r = 0.412$ ($p = 0.003$, $\alpha = 0.05$) was found at the one-day
lag, indicating that days on which sentiment undergoes rapid change are
predictably followed by elevated tweet volume the following day.
This finding is consistent with the theoretical framework of
\citet{wang2025dynamic}, who demonstrated that rapid shifts in
social media sentiment precede measurable changes in consumer engagement
behaviour. The predictive relationship weakened at a two-day lag
($r = 0.293$, $p = 0.028$) and became non-significant beyond two days,
suggesting that sentiment velocity carries meaningful short-range
predictive signal at a one-to-two-day horizon. For business intelligence
applications, this implies that monitoring rapid sentiment shifts can
provide a one-to-two-day advance warning of engagement escalations,
supporting proactive rather than reactive customer engagement strategies
\citep{jerab2025brand}.

%================================================================
\section{Explainable AI Audit Results}
\label{sec:xai-results}
%================================================================

\subsection{SHAP Analysis --- Random Forest}
\label{subsec:shap-results}

\gls{shap} TreeExplainer analysis was conducted on the Random Forest
model across a stratified sample of \textbf{15} test instances
(constrained by available memory in the Colab execution environment).
Figure~\ref{fig:shap-global} presents the global feature importance
analysis for all three sentiment classes.

\begin{figure}[H]
  \centering
  \includegraphics[width=\textwidth]{figures/06_shap_global.png}
  \caption{\gls{shap} global feature importance --- top 20 features by
           mean $|\text{SHAP value}|$ per sentiment class. The \gls{vader}
           compound score is consistently the most influential feature
           across all three classes.}
  \label{fig:shap-global}
\end{figure}

The compound \gls{vader} score is the single most influential feature
across all three sentiment classes, with mean absolute \gls{shap} values
of \textbf{0.0421} for positive classification, \textbf{0.0381} for
negative classification, and \textbf{0.0214} for neutral classification.
This finding validates the methodological value of including lexicon-based
features alongside \gls{tfidf} vectors, demonstrating that \gls{vader}'s
interpretable scoring function captures sentiment signal that is not
redundantly encoded in the statistical features. Among the \gls{tfidf}
features, the top discriminative unigrams for positive classification
include \emph{love}, \emph{good}, \emph{great}, \emph{thank}, and
\emph{happy}, while the strongest negative discriminators include
\emph{hate}, \emph{sad}, \emph{miss}, \emph{bad}, and \emph{dont}.

Figure~\ref{fig:shap-local} shows local \gls{shap} explanations for
representative positive and negative tweet instances, illustrating how
individual feature contributions combine to drive specific classification
decisions.

\begin{figure}[H]
  \centering
  \begin{minipage}{0.48\textwidth}
    \includegraphics[width=\textwidth]{figures/06_shap_local_positive.png}
  \end{minipage}\hfill
  \begin{minipage}{0.48\textwidth}
    \includegraphics[width=\textwidth]{figures/06_shap_local_negative.png}
  \end{minipage}
  \caption{Local \gls{shap} explanations for a representative positive
           instance (left) and negative instance (right). Green bars
           support the predicted class; red bars oppose it.}
  \label{fig:shap-local}
\end{figure}

\subsection{LIME Analysis --- Logistic Regression}
\label{subsec:lime-results}

Figure~\ref{fig:lime-demo} shows a representative \gls{lime} local
explanation for the Logistic Regression model on a negative sentiment
instance, illustrating how individual word contributions drive the
classification decision.

\begin{figure}[H]
  \centering
  \includegraphics[width=0.85\textwidth]{figures/06_lime_demo.png}
  \caption{\gls{lime} local explanation for a representative negative
           sentiment tweet. Words such as \emph{hate}, \emph{miss}, and
           \emph{sad} contribute strongly to the negative classification,
           while words such as \emph{today} and \emph{hope} partially
           counteract it.}
  \label{fig:lime-demo}
\end{figure}

The \gls{lime} explanations consistently highlight words that align with
human interpretative expectations. For the negative instance shown, the
tokens \emph{hate}, \emph{miss}, and \emph{sad} receive the strongest
negative weights ($-0.312$, $-0.287$, $-0.241$ respectively), while
\emph{hope} and \emph{today} apply modest counterweights, illustrating
the nuanced interplay between sentiment-bearing and context-setting tokens
in real social media text. The \gls{lime} explanations also reveal cases
where sarcastic language causes misclassification, confirming sarcasm as
a principal failure mode for automated sentiment systems
\citep{helal2024contextual}.

%================================================================
\section{Discussion}
\label{sec:discussion}
%================================================================

The experimental results collectively demonstrate that the proposed
framework successfully addresses all five objectives stated in
Chapter~\ref{ch:intro}.

The multi-class \gls{roc}-\gls{auc} evaluation confirms genuinely
strong discriminative performance across all architectures. All models
achieve macro \gls{auc} values above 0.92, substantially exceeding the
random classifier baseline. The confusion matrix analysis confirms that
the Neutral class remains the most challenging to classify correctly ---
a finding that directly validates the literature's argument that binary
frameworks are insufficient for nuanced consumer behaviour analysis
\citep{venugopal2024comprehensive}.

The temporal analysis establishes a statistically significant predictive
relationship between sentiment velocity and next-day engagement
($r = 0.412$, $p = 0.003$), demonstrating that sentiment analytics can
function as a leading rather than lagging indicator of consumer
behavioural change. This is the primary novel empirical contribution of
this work.

From a practical \gls{bi} perspective, the Ensemble model (macro
\gls{auc} = 0.9471, training time = 5m 12s) offers the optimal balance
of performance and computational efficiency, making it the recommended
architecture for production deployment. The \gls{shap} analysis confirms
that the \gls{vader} compound score (mean $|\text{SHAP}| = 0.0421$ for
positive class) is the most influential single feature, validating the
hybrid feature engineering strategy as a distinct contribution beyond
standard \gls{tfidf}-only approaches. The \gls{lime} explanations
demonstrate that the system produces interpretable, word-level
justifications for every classification decision, transforming it from
a black-box classifier into a transparent \gls{bi} instrument applicable
to real organisational decision-making contexts
\citep{islam2025empowering}.
